

\chapter{Conditional expectations}

\section{Elementary conditional expectations}
Let $X$ and $Y$ be random variables defined on the
same probability space
$(\Omega, \mathcal{F}, \mathrm{P})$. We wish to
generalize the expectation operator $\mathbf{E}$ that
takes a random variable $X$ and returns a number, to
include the possibility of having additional
information, encoded by the random variable $Y$.
Thus, we want to define the
$\textbf{conditional expectation}$$\mathbf{E}\left(X\,|\,Y\right)$ (``the conditional
expectation of $X$ given $Y$''). This is tricky to do
for general random variables, so let's start with the
simple case in which $X$ and $Y$ are both discrete
random variables taking on finitely many values with
positive probability. That is, assume there are
numbers $x_1,\ldots, x_m$ and $y_1,\ldots, y_n$ such
that

\[
\mathrm{P}\left( X \in \{ x_1,\ldots, x_m\}, Y \in \{ y_1,\ldots, y_n \} \right) = 1.
\]

\noindent In this case, if we know that $Y=y_j$ for some
$1\le j\le n$, then the conditional distribution of
$X$ becomes

\[
\mathrm{P}(X = x_i\ |\ Y=y_j) = \frac{\mathrm{P}(X=x_i, Y=y_j)}{\mathrm{P}(Y=y_j)}, \qquad 1\le i\le m.
\]

\noindent The expected value of this distribution is

\[
\mathbf{E}\left(X\,|\,Y=y_j\right) = \sum_{i=1}^n \mathrm{P}(X = x_i\ |\ Y=y_j) x_i.
\]

\noindent For each $j$, this is a \emph{number}. It makes sense to
define a \emph{random variable}, not a number, to be the
random variable $Z=\mathbf{E}\left(X\,|\,Y\right)$
that is equal to
$\mathbf{E}\left(X=x_i\,|\,Y=y_j\right)$ on the event
$\{Y=y_j\}$:

\[
Z = \sum_{j=1}^n \mathbf{E}\left(X\,|\,Y=y_j\right) \mathbf{1}_{\{Y=y_j\}}.
\]

\noindent Let us try to think of a more conceptual way to look
at this definition. First of all, note that the actual
values $y_1,\ldots, y_n$ that the random variable $Y$
takes are not so important. Rather, what is important
is that $Y$ partitions the probability space $\Omega$
into disjoint subsets, such that knowing which of the
subsets we end up in changes our view of the
distribution of $X$. Therefore in the definition of
$Z$ we could insert in place of $Y$ an object that
encodes this slightly lesser information (i.e., not
including the values of $Y$). It turns out that the
correct object to use is the $\sigma$-algebra
$\sigma(Y)$ generated by $Y$. For a general r.v.,
this is defined by

\[
\sigma(Y) = \big\{ \{Y \in B\}\ :\ B\in\mathcal{B}(\R) \big\}
= \{ Y^{-1}(B)\ :\ B\in\mathcal{B}(\R) \}
\]

\noindent and is a sub-$\sigma$-algebra of $\mathcal{F}$. For
our discrete $Y$, it is the $\sigma$-algebra
generated by the events $\{Y=y_j\}, \ j=1,\ldots,n$.
So, given a sub-$\sigma$-algebra
$\mathcal{G} \subset \mathcal{F}$ generated by a
family of disjoint subsets
$A_1,\ldots,A_n \in \mathcal{F}$ that partitions
$\Omega$, we can rewrite the definition of the
conditional expectation as the random variable
$Z = \mathbf{E}\left(X\,|\,\mathcal{G}\right)$ (``the
conditional expectation of $X$ given the
$\sigma$-algebra $\mathcal{G}$'') given by

\[
Z = \sum_{j=1}^n \mathbf{E}\left(X\,|\,A_j\right) \mathbf{1}_{A_j}.
\]

\noindent Here, for each $1\le j\le n$ the quantity
$\mathbf{E}\left(X\,|\,A_j\right)$ is a number, the
expected value of the conditional distribution of $X$
on the event $A_j$.

The next observation is that $Z$ satisfies two
properties, which seem rather trivial and
uninteresting in this context, but will turn out to be
crucial to generalizing the definition of
$\mathbf{E}\left(X\,|\,Y\right)$ to arbitrary random
variables:

\begin{enumerate}[label=(\roman*)]
\item $Z$ is measurable with respect to $\mathcal{G}$. That
is, for each Borel set $B\in \mathcal{B}(\R)$, the
event $\{Z\in B\} \in \mathcal{G}$. Equivalent ways
of saying the same thing are that
$\sigma(Z)\subset \mathcal{G}$, or that $Z$ would
remain a random variable even if we change the
probability space to
$(\Omega,\mathcal{G},\mathrm{P})$.

\item For any event $E \in \mathcal{G}$, we have
$\mathbf{E}(Z \mathbf{1}_E) = \mathbf{E}(X
\mathbf{1}_E)$.
\end{enumerate}

\begin{xexercise}
Prove the above two properties.
\end{xexercise}

\section{General conditional expectations}
The above discussion brings us to the following
important definition.

\begin{defn}
Let $X$ be a random variable on
$(\Omega, \mathcal{F}, \mathrm{P})$ with
$\mathbf{E}|X|<\infty$, and let $\mathcal{G}$ be a
sub-$\sigma$-algebra of $\mathcal{F}$. If a random
variable $Z$ satisfies properties (i) and (ii) above,
we say that $Z$ is the \textbf{conditional expectation of $X$
given $\mathcal{G}$}, and denote
$Z = \mathbf{E}\left(X\,|\,\mathcal{G}\right)$. If $Y$
is another random variable on
$(\Omega, \mathcal{F}, \mathrm{P})$, then taking
$\mathcal{G}=\sigma(Y)$ we may denote
$Z = \mathbf{E}\left(X\,|\,Y\right) =
\mathbf{E}\left(X\,|\,\sigma(Y)\right)$,
and refer to $Z$ as the \textbf{conditional expectation of
$X$ given $Y$}.
\end{defn}

\noindent The use of the notation
$Z = \mathbf{E}\left(X\,|\,\mathcal{G}\right)$
contains the implicit assumption that the conditional
expectation is unique. This is true with the
convention that random variables are defined only up
to almost sure equivalence. That is, if $Z=Z'$ a.s.\
we consider $Z$ and $Z'$ to be the same random
variable (in particular, any event involving $Z$ will
have the same probability as the corresponding event
involving $Z'$). In fact, we will prove the following
result which also answers the key questions regarding
existence and integrability.

\begin{thm}
\label{thm-cond-ext}
The conditional expectation exists, is an integrable
random variable, and is unique up to almost sure
equivalence.
\end{thm}

\begin{proof}[Proof of integrability]
Let $A = \{ Z>0 \}$. By (i), $A, A^c \in \mathcal{G}$.
Therefore by (ii), we have

\begin{align*}
\mathbf{E}(Z \mathbf{1}_A) &= \mathbf{E}(X \mathbf{1}_A) \le \mathbf{E}\left( |X| \mathbf{1}_A \right), \\
\mathbf{E}(Z \mathbf{1}_{A^c}) &= \mathbf{E}(X \mathbf{1}_{A^c}) \ge -\mathbf{E}\left( |X| \mathbf{1}_{A^c} \right),
\end{align*}

\noindent Subtracting the two equations gives

\[
\mathbf{E}|Z| = \mathbf{E}(Z\mathbf{1}_A - Z\mathbf{1}_{A^c}) \le \mathbf{E}\left( |X| (\mathbf{1}_A + \mathbf{1}_{A^c}) \right) = \mathbf{E}|X| < \infty.
\]

\end{proof}

\begin{proof}[Proof of uniqueness]
Assume that $Z$ and $Z'$ both satisfy properties (i),
(ii). It follows that

\[
\mathbf{E}(Z \mathbf{1}_A) = \mathbf{E}(Z' \mathbf{1}_A)\ \  \textrm{ for all }A \in \mathcal{G}.
\]

\noindent For some $\epsilon>0$, take
$A = \{ Z-Z' \ge \epsilon \}$ to get that

\[
0 = \mathbf{E}( (Z-Z')\mathbf{1}_A ) \ge \epsilon \mathrm{P}(A),
\]

\noindent so $\mathrm{P}(A) = 0$. Since $\epsilon$ was an
arbitrary positive number, this implies that $Z\le Z'$
a.s. Reversing the roles of $Z$ and $Z'$ gives that
also $Z'\le Z$ a.s., so $Z=Z'$ a.s., which proves the
uniqueness claim.
\end{proof}

\noindent The most tricky part in defining conditional
expectations is proving that they exist. The usual
proof involves an application of an important result
from measure theory, the Radon-Nikodym theorem, which
we review in the next section.

\section{Absolute continuity and the Radon-Nikodym theorem}
We will formulate a version of the Radon-Nikodym
theorem for probability spaces. The theorem involves
the useful concept of absolute continuity of measures.

\begin{defn}
Let $P$ and $Q$ be two probability measures on a
measurable space $(\Omega, \mathcal{F})$. We say that
\textbf{$P$ is absolutely continuous with respect to $Q$},
and denote $P<<Q$, if for any $A\in\mathcal{F}$, if
$Q(A)=0$ then $P(A)=0$.
\end{defn}

\noindent There is an interesting intuitive interpretation to
the notion of absolute continuity. In many real-life
situations involving probability we know the
measurable space $(\Omega, \mathcal{F})$ but do not
know which probability measure correctly describes the
statistical distribution of the outcome of the
experiment (for example, someone hands us a die to
roll but we are not sure if it is a fair or loaded
die). Say there are two possible measures, $P$ and
$Q$, which we are considering. It makes sense to
perform the experiment and try to guess which of $P$
and $Q$ is the correct one based on the result. If we
are lucky, the result might fall into an event
$A \in \mathcal{F}$ which has the property that
$Q(A)=0$ but $P(A)>0$, in which case we will know
that we can safely rule out $Q$. The relation $P<<Q$
means that this cannot happen; i.e., based on a single
experiment, or even a finite number of repeated
experiments, we can never rule out $Q$ as the correct
measure, although we may become increasingly convinced
that $P$ is the correct one as the statistical
evidence mounts\footnote{The question of precisely how fast can one become
convinced of the correct answer leads to the concept
of $\textbf{relative entropy}$, which will be the
subject of a future homework problem.}. If $P<<Q$
and $Q<<P$, meaning both measures are mutually
absolutely continuous w.r.t.\ each other, we can also
never rule out $P$ so we may never be entirely sure
which measure correctly describes the experiment.

The following lemma sheds further light on the notion
of absolute continuity. Its proof is a nice
application of the first Borel-Cantelli lemma.

\begin{lem}
\label{lem-abs-cont}
$P$ is absolutely continuous with respect to $Q$ if
and only if for any $\epsilon>0$ there exists a
$\delta>0$ such that if $Q(A)<\delta$ then
$P(A)<\epsilon$.
\end{lem}

\begin{proof}
The ``if'' direction is immediate. To prove the
``only if,'' assume the negation of the condition
``for any $\epsilon>0$, $\ldots$'' That is, there
must exist an $\epsilon>0$ such that for any
$\delta>0$ there is an event $A_\delta$ such that
$Q(A_\delta)< \delta$ but $P(A_\delta)\ge\epsilon$.
Now consider the event

\[
E= \{ A_{1/{2^n}}\textrm{ i.o.} \} =  \bigcap_{N=1}^\infty \bigcup_{n=N}^\infty A_{1/2^n}
\]

\noindent Since $\sum_{n=1}^\infty 2^{-n}<\infty$, by the
Borel-Cantelli lemma we get that $Q(E)=0$. On the
other hand,

\[
P(E) = \lim_{N\to\infty} P\left[\bigcup_{n=N}^\infty A_{1/2^n}\right] \ge \limsup_{N\to\infty} P(A_{1/2^n}) \ge \epsilon>0.
\]

\noindent This shows that $P$ is not absolutely continuous with
respect to $Q$.
\end{proof}

\noindent If $Q$ is a probability measure on
$(\Omega, \mathcal{F})$, one way to construct a
probability measure $P$ for which $P<<Q$ is to let

\begin{equation*} \tag{5}\label{eq:abs-cont-p-q}
P(E) = \mathbf{E}_{Q}(Z \mathbf{1}_E), \qquad (E\in \mathcal{F}).
\end{equation*}

\noindent where $Z$ is some nonnegative random variable on the
probability space $(\Omega,\mathcal{F},Q)$ satisfying
$\mathbf{E}(Z)=1$ (the notation
$\mathbf{E}_{Q}(\cdot)$ emphasizes that the
expectation operator is the one associated with the
probability measure $Q$). The Radon-Nikodym theorem
says that this construction is the most general one
possible.

\begin{thm}[The Radon-Nikodym theorem]
$P<<Q$ if and only if there exists a random variable
$Z$ on $(\Omega, \mathcal{F}, \mathrm{P})$, unique up
to $Q$-a.s.-equivalence, such that $Z\ge 0$ a.s.\ and
the relation $P(E) = \mathbf{E}_{Q} (Z \mathbf{1}_E)$
holds for any $E\in \mathcal{F}$. The random variable
$Z$ is referred to as the
$\textbf{Radon-Nikodym derivative}$ of $P$ relative
to $Q$ and denoted

\[
Z = \frac{dP}{dQ}.
\]

\end{thm}

\noindent As mentioned above, the ``if'' part of the theorem is
immediate, so it is really the ``only if'' part which
is interesting. The Radon-Nikodym theorem has several
classical proofs which may be read in measure theory
and analysis textbooks. We will assume it without
proof for now, and later sketch a probabilistic proof
that uses martingale theory. (This is almost circular
logic, since we will use conditional expectations to
develop martingales and conditional expectations rely
on the Radon-Nikodym theorem. However, with a bit of
care one can make sure to use only the ``elementary''
version of conditional expectations and thus obtain a
genuinely new proof that does not make use of circular
reasoning --- this is one of the nice examples of
probability theory ``giving back'' to the rest of
mathematics.)

Note that the Radon-Nikodym derivative generalizes
the ordinary derivative: if $Q$ is Lebesgue measure
on the measure space $((0,1),\mathcal{B})$ and $P<<Q$,
then the random variable $Z$ in (\ref{eq:abs-cont-p-q})
can be thought of as the density function of a random
variable whose distribution measure is $P$, and can
be computed as an actual derivative of the cumulative
distribution function of this random variable, i.e.,

\[
Z(x) = \frac{d}{dx} P((0,x)) \quad\textrm{a.s.}
\]

\noindent Another reason why it makes sense to think of $Z$ as
a kind of ``derivative'' is that if one uses the
Lebesgue integral notation $\int_E Z dQ$ instead of
our more probabilistic notation
$\mathbf{E}_{Q}(Z\mathbf{1}_E)$, the relation
(\ref{eq:abs-cont-p-q}) can be rewritten as an intuitive
``change of measures'' formula

\begin{equation*} \tag{6}\label{eq:abs-cont-change-var}
\int_E dP = \int_E \frac{dP}{dQ} dQ.
\end{equation*}

\begin{xexercise}
Prove that if $P<<Q$ are two probability measures on
a measurable space $(\Omega, \mathcal{F})$ then
(\ref{eq:abs-cont-change-var}) generalizes to the identity

\[
\int_E X dP = \int_E X \frac{dP}{dQ} dQ
\]

\noindent which holds for any random variable $X$ for which
$\mathbf{E}_P|X|<\infty$.
\end{xexercise}

\begin{proof}[Proof of existence in Theorem \ref{thm-cond-ext}]
We are now in a position to prove the existence of
conditional expectations. Assume first that $X\ge 0$.
Define the measure $Q$ on the measurable space
$(\Omega, \mathcal{G})$ by

\[
Q(E) = \mathbf{E}_{\mathrm{P}}(X \mathbf{1}_E), \qquad E\in\mathcal{G}.
\]

\noindent Note that $Q$ is not a \emph{probability} measure; rather,
it is a \emph{finite} measure, which is like a probability
measure but takes values in $[0,\infty)$ and is not
required to satisfy $Q(\Omega)=1$. However, the
Radon-Nikodym theorem remains true for such measures
(even more generally for so-called \emph{$\sigma$-finite}
measures), as can easily be seen by replacing
$Q(\cdot)$ by its scalar multiple
$Q(\Omega)^{-1} Q(\cdot)$, which is a probability
measure. Now, the original probability measure
$\mathrm{P}$ on $(\Omega, \mathcal{F})$ can also be
thought of as a measure on the measurable space
$(\Omega, \mathcal{G})$. Furthermore, it is easy to
see that $Q << \mathrm{P}$. Thus, we are exactly in
the situation described in the Radon-Nikodym theorem,
and we conclude that there is a random variable
$Z=dQ/dP$, measurable with respect to the
$\sigma$-algebra $\mathcal{G}$, such that for any
$E\in\mathcal{G}$ we have

\[
\mathbf{E}_{\mathrm{P}}(Z \mathbf{1}_E) = Q(E) = \mathbf{E}_{\mathrm{P}}(X \mathbf{1}_E).
\]

\noindent Thus, $Z$ satisfies the two properties (i)--(ii) in
the definition of the conditional expectation, so the
conditional expectation
$Z=\mathbf{E}\left(X\,|\,\mathcal{G}\right)$ exists.

Finally, for a general random variable $X$ with
$\mathbf{E}|X|<\infty$, write $X=X_+-X_-$ where
$X_+, X_-\ge 0$ are the non-negative and non-positive
parts of $X$, respectively. It is immediate to check
that the random variable
$Z = \mathbf{E}\left(X_+\,|\,\mathcal{G}\right) -
\mathbf{E}\left(X_-\,|\,\mathcal{G}\right)$
satisfies the properties (i)--(ii) required of the
conditional expectation
$\mathbf{E}\left(X\,|\,\mathcal{G}\right)$, so again
we have shown that the conditional expectation exists.
\end{proof}

\section{Properties and examples of conditional expectations}

\begin{enumerate}[label=\arabic*.]
\item If $X$ is measurable with respect to $\mathcal{G}$
(i.e., $\sigma(X) \subset \mathcal{G}$) then
$\mathbf{E}\left(X\,|\,\mathcal{G}\right) = X$. The
intuition is that if the information given to us
(encoded by the $\sigma$-algebra $\mathcal{G}$) is
enough to tell the value of $X$ (see the exercise
below), then our best guess for the average value of
$X$ is $X$ itself.

\begin{xexercise}
Let $X$ and $Y$ be random variables defined on the
same probability space
$(\Omega, \mathcal{F}, \mathrm{P})$. Show that
$\sigma(X) \subset \sigma(Y)$ (in this case we say
that \textbf{$X$ is measurable with respect to $Y$}) if and
only if there exists a Borel-measurable function
$h:\R\to\R$ such that $X=h(Y)$ a.s.
\end{xexercise}

\item In particular, if $\mathcal{G}=\mathcal{F}$ is the
original $\sigma$-algebra of the probability space
$(\Omega, \mathcal{F}, \mathrm{P})$, then
$\mathbf{E}\left(X\,|\,\mathcal{G}\right) = X$.

\item At the opposite extreme, if
$\mathcal{G}= \{ \emptyset, \Omega \}$, then
$\mathbf{E}\left(X\,|\,\mathcal{G}\right) =
\mathbf{E}(X)$,
the usual expectation of $X$. Intuitively, the
minimal $\sigma$-algebra $\{\emptyset, \Omega\}$
represents zero information.

\item More generally, if $X$ is independent of $\mathcal{G}$
then
$\mathbf{E}\left(X\,|\,\mathcal{G}\right)=\mathbf{E}(X)$,
since $\mathbf{E}(X)$ is a $\mathcal{G}$-measurable
random variable and for any $E\in\mathcal{G}$, since
$X$ is independent of $\mathbf{1}_E$,

\[
\mathbf{E}(X\mathbf{1}_E) = \mathbf{E}(X)\mathbf{E}(\mathbf{1}_E) = \mathbf{E}(\mathbf{E}(X)\mathbf{1}_E).
\]

\item Conditional expectation is a linear operator:
$\mathbf{E}\left(aX+bY\,|\,\mathcal{G}\right) = a
\mathbf{E}\left(X\,|\,\mathcal{G}\right) + b
\mathbf{E}\left(Y\,|\,\mathcal{G}\right).$

\begin{proof}
The random variable
$Z=a \mathbf{E}\left(X\,|\,\mathcal{G}\right) + b
\mathbf{E}\left(Y\,|\,\mathcal{G}\right)$
is $\mathcal{G}$-measurable and for any
$E\in\mathcal{G}$ we have

\begin{align*}
\mathbf{E}(Z\mathbf{1}_E) &= a\mathbf{E}(\mathbf{E}\left(X\,|\,\mathcal{G}\right)\mathbf{1}_E)+b\mathbf{E}(\mathbf{E}\left(Y\,|\,\mathcal{G}\right)\mathbf{1}_E) =
a \mathbf{E}(X\mathbf{1}_E)+b\mathbf{E}(Y\mathbf{1}_E) \\&=\mathbf{E}((aX+bY)\mathbf{1}_E),
\end{align*}

\noindent so $Z$ satisfies the requirements that qualify it to
be the conditional expectation of $aX+bY$.
\end{proof}

\item Monotonicity: If $X\le Y$ then
$\mathbf{E}\left(X\,|\,\mathcal{G}\right)\le\mathbf{E}\left(Y\,|\,\mathcal{G}\right)$
a.s.

\begin{proof}
For any $E\in\mathcal{G}$, we have that

\[
\mathbf{E}(\mathbf{E}\left(X\,|\,\mathcal{G}\right)\mathbf{1}_E) = \mathbf{E}(X\mathbf{1}_E) \le
\mathbf{E}(Y\mathbf{1}_E) \le \mathbf{E}(\mathbf{E}\left(Y\,|\,\mathcal{G}\right)\mathbf{1}_E).
\]

\noindent Taking
$E=\{\mathbf{E}\left(X\,|\,\mathcal{G}\right)-\mathbf{E}\left(Y\,|\,\mathcal{G}\right)\ge
\epsilon\}$
for some arbitrary $\epsilon>0$, we get that $E$ must
have probability $0$. Since this is true for any
$\epsilon>0$, the result follows.
\end{proof}

\item Conditional form of the monotone convergence theorem:
If $X_n\ge 0$ and $X_n \uparrow X$ a.s.\ as
$n\to\infty$, where $\mathbf{E}(X)<\infty$, then
$\mathbf{E}\left(X_n\,|\,\mathcal{G}\right)\nearrow
\mathbf{E}\left(X\,|\,\mathcal{G}\right)$
a.s.\ as $n\to\infty$.

\begin{proof}
By the monotonicity proved above, the conditional
expectations
$\mathbf{E}\left(X_n\,|\,\mathcal{G}\right)$ are
increasing to an a.s.\ limiting r.v.\ $Z$, which is
$\mathcal{G}$-measurable. Note that
$\mathbf{E}\left(X_n\,|\,\mathcal{G}\right)\le
\mathbf{E}\left(X\,|\,\mathcal{G}\right)$,
which is an integrable upper bound, so we can apply
the dominated convergence theorem and deduce that for
every $E\in \mathcal{G}$,

\[
\mathbf{E}(Z\mathbf{1}_E) = \mathbf{E}(\lim_{n\to\infty} \mathbf{E}\left(X_n\,|\,\mathcal{G}\right)\mathbf{1}_E)
= \lim_{n\to\infty} \mathbf{E}(\mathbf{E}\left(X_n\,|\,\mathcal{G}\right)\mathbf{1}_E) = \lim_{n\to\infty} \mathbf{E}(X_n\mathbf{1}_E)
= \mathbf{E}(X\mathbf{1}_E).
\]

\noindent Thus $Z$, qualifies as the conditional expectation of
$X$ given $\mathcal{G}$.
\end{proof}

\item Conditional form of Jensen's inequality: If
$\varphi:\R\to\R$ is a convex function and
$\mathbf{E}|X|, \mathbf{E}|\varphi(X)|<\infty$ then
$\varphi(\mathbf{E}\left(X\,|\,\mathcal{G}\right))
\le
\mathbf{E}\left(\varphi(X)\,|\,\mathcal{G}\right).$

\begin{proof}
The proof is similar to the non-conditional case,
with a small twist: we observe that $\varphi(x)$ is a
supremum of a countable number of linear functions,
namely

\[
\varphi(x) = \sup \{ ax+b\,:\,a,b\in \mathbb{Q}, ax+b\le \varphi(x)\textrm{ for all }x\},
\]

\noindent then continue as in the original proof. The
restriction to a countable number of linear functions
is necessary since all statements pertaining to
conditional expectations are almost sure statements,
which means that there is an exceptional
$0$-probability set being discarded, and we can only
allow a countable number of these.
\end{proof}

\item For $p\ge 1$,
$\mathbf{E}\left(
\left|\mathbf{E}\left(X\,|\,\mathcal{G}\right)\right|^p
\right) \le \mathbf{E}|X|^p$
(i.e., the conditional expectation operator is a
contraction in $L_p(\Omega,\mathcal{G},\mathrm{P})$.)

\begin{proof}
By Jensen's inequality,
$|\mathbf{E}\left(X\,|\,\mathcal{G}\right)|^p \le
\mathbf{E}\left(|X|^p\,|\,\mathcal{G}\right)$.
Now take expectations and use the fact that:
\end{proof}

\item $\mathbf{E}(\mathbf{E}\left(X\,|\,\mathcal{G}\right))
= \mathbf{E}(X)$.
This is a special case of:

\item If
$\mathcal{G}_1\subseteq \mathcal{G}_2 \subseteq
\mathcal{F}$
are sub-$\sigma$-algebras then we have

\begin{align*}
\mathbf{E}\left(\mathbf{E}\left(X\,|\,\mathcal{G}_2\right)\,|\,\mathcal{G}_1\right) &= \mathbf{E}\left(X\,|\,\mathcal{G}_1\right), \\
\mathbf{E}\left(\mathbf{E}\left(X\,|\,\mathcal{G}_1\right)\,|\,\mathcal{G}_2\right) &= \mathbf{E}\left(X\,|\,\mathcal{G}_1\right),
\qquad  \textrm{(no, there is no typo here\ldots)}
\end{align*}

\begin{xexercise}
Prove this.
\end{xexercise}

\item If $X$ is $\mathcal{G}$-measurable, then for any
random variable $Y$, if
$\mathbf{E}|Y|, \mathbf{E}|XY|<\infty$ then

\[
\mathbf{E}\left(XY\,|\,\mathcal{G}\right) = X \mathbf{E}\left(Y\,|\,\mathcal{G}\right).
\]

\noindent The idea here is that if $X$ is measurable with
respect to $\mathcal{G}$ then it appears like a
constant from the point of view of conditional
expectations given the information contained in
$\mathcal{G}$ and can therefore be ``pulled outside''
of the expectation.

\begin{proof}
The random variable
$Z=X\mathbf{E}\left(Y\,|\,\mathcal{G}\right)$ is
$\mathcal{G}$-measurable, so we need to check that
for any $E\in \mathcal{G}$,

\begin{equation*} \tag{7}\label{eq:condexp-x-times-y}
\mathbf{E}(Z\mathbf{1}_E) = \mathbf{E}(XY \mathbf{1}_E).
\end{equation*}

\noindent First, if $X=\mathbf{1}_B$ for some $B\in\mathcal{G}$
then

\[
\mathbf{E}(Z \mathbf{1}_E) = \mathbf{E}(\mathbf{E}\left(Y\,|\,\mathcal{G}\right)\mathbf{1}_{B\cap E}) = \mathbf{E}(Y\mathbf{1}_{B\cap E})
= \mathbf{E}(XY \mathbf{1}_E)
\]

\noindent so the claim is true. By linearity it follows that if
$X$ is a ``simple'' random variable (a linear
combination of finitely many indicator variables,
i.e., a random variable taking finitely many values)
then (\ref{eq:condexp-x-times-y}) still holds. Now, assume
that $X,Y\ge 0$, then we can take a sequence
$(X_n)_{n=1}^\infty$ of simple random variables that
are $\mathcal{G}$-measurable and such that
$X_n\uparrow X$. By the monotone converegence
theorem, it follows that (\ref{eq:condexp-x-times-y})
holds also in this case. Finally, for general $X,Y$
one proves (\ref{eq:condexp-x-times-y}) by splitting $X$
and $Y$ into their positive and negative parts.
\end{proof}
\end{enumerate}

\section{Conditional expectation as the least mean square error estimator}
In this section we show that the conditional
expectation is actually the solution to a very natural
estimation problem. Let $X$ be a random variable with
finite variance. Assume that we wish to give the best
possible estimate for the value of $X$, but only have
access to the information encoded by the
$\sigma$-algebra $\mathcal{G}$ (i.e., for each event
$E\in\mathcal{G}$ we know if $E$ occurred or did not
occur). Our estimate will therefore be some random
variable $Y$ which is $\mathcal{G}$-measurable. A
natural measure for the quality of the estimate is the
$\textbf{mean square error}$

\[
\operatorname{MSE}_X(Y) = \mathbf{E}(X-Y)^2.
\]

\noindent The problem of finding the $Y$ that minimizes this
error is very natural and of great practical
importance. Its solution is given in the following
theorem.

\begin{thm}
If $\mathbf{E}(X^2)<\infty$, then
$\mathbf{E}\left(X\,|\,\mathcal{G}\right)$ is the
unique (up to a.s.\ equivalence) random variable $Y$
that minimizes the mean square error
$\operatorname{MSE}_X(Y)$ among all
$\mathcal{G}$-measurable random variables.
\end{thm}

\begin{proof}
Denote $Z=\mathbf{E}\left(X\,|\,\mathcal{G}\right)$.
If $Y$ is a $\mathcal{G}$-measurable random variable
with $\mathbf{E} Y^2<\infty$ (if $Y$ does not have
finite variance then clearly the MSE will be
infinite), then

\begin{align*}
\mathbf{E}(X-Y)^2 &= \mathbf{E}[((X-Z)-(Y-Z))^2] \\
&= \mathbf{E}(X-Z)^2 + \mathbf{E}(Y-Z)^2
+ 2 \mathbf{E}[(X-Z)(Y-Z)]
\end{align*}

\noindent Denote $W=Y-Z$. This is a $\mathcal{G}$-measurable
random variable, so by one of the properties proved in
the previous section,

\[
\mathbf{E}(WX) = \mathbf{E}(\mathbf{E}\left(WX\,|\,\mathcal{G}\right))
= \mathbf{E}(W \mathbf{E}\left(X\,|\,\mathcal{G}\right)) = \mathbf{E}(WZ)
\]

\noindent (note that
$\mathbf{E}|WX|\le (\mathbf{E} W^2 \mathbf{E}
X^2)^{1/2}<\infty$
by the Cauchy-Schwartz inequality) and therefore
$\mathbf{E}[(X-Z)(Y-Z)] = \mathbf{E}[(X-Z)W]=0$. We
get that

\[
\mathbf{E}(X-Y)^2 = \mathbf{E}(X-Z)^2 + \mathbf{E}(Y-Z)^2,
\]

\noindent which is clearly minimized when (and only when) $Y=Z$
a.s.
\end{proof}

\noindent The above result has an interesting geometric
interpretation. The expression $\mathbf{E}(X-Y)^2$ is
the square of the $L_2$-distance between $X$ and $Y$
in the Hilbert space
$L_2(\Omega, \mathcal{F},\mathrm{P})$. The theorem
says that the conditional expectation operator
$\mathbf{E}\left(\cdot\,|\,\mathcal{G}\right)$ takes
a random variable $X$ and returns the closest point
to $X$ in the linear subspace
$L_2(\Omega, \mathcal{G},\mathrm{P})\subset
L_2(\Omega, \mathcal{F},\mathrm{P})$
consisting of square-integrable
$\mathcal{G}$-measurable random variables. By
elementary properties of Hilbert spaces, this is
equivalent to the statement that
$\mathbf{E}\left(\cdot\,|\,\mathcal{G}\right)$ is the
orthogonal projection operator onto the subspace
$L_2(\Omega, \mathcal{G},\mathrm{P})$.

One way in which this geometric interpretation is
useful is that it suggests an alternative approach to
defining conditional expectations that does not rely
on the Radon-Nikodym Theorem: first construct the
conditional expectation operator
$\mathbf{E}\left(\cdot\,|\,\mathcal{G}\right)$ for
square-integrable random variables by defining it as
an orthogonal projection operator; then extend the
definition to the space
$L_1(\Omega,\mathcal{F},\mathrm{P})$ of integrable
random variables by approximating such variables by
square-integrable ones. The book \emph{Probability With
Martingales} by David Williams is one textbook where
such an approach is developed.
